% GENERATED FILE. Edit canonical lessons, then run npm run generate:books.
% canonical-source-hash: fnv1a64:ad91950d59a0a5d4
% canonical-lessons: ES-C541-fase, ES-C541-transcurso, ES-C541-duracion, ES-C541-prolongarse, ES-C541-transcurrir

\chapter{Cuánto duró}
\label{ch:es-cuanto-duro}
\hlchaptermodality{541}

\begin{chapteropening}
\textbf{By the end of this chapter:} \emph{I can say how long something went on, what phase it was in and how the time went by.}

Complete the last lesson of chapter 541: I can say how long something went on, what phase it was in and how the time went by.
\end{chapteropening}

\section[la fase]{la fase — the phase, the stage}

\label{lesson:ES-C541-fase}



\begin{quote}
Before the new one: say the Spanish for weekly, then the Spanish for yearly.
\end{quote}

\subsection*{You'll want to know: la fase}
\textbf{la fase} — \textquotedblleft{}the phase, the stage\textquotedblright{}.

\begin{quote}
El proyecto está en la primera fase.
\end{quote}

\begin{quote}
\emph{The project is in its first phase.}
\end{quote}

\begin{cousinweb}
\emph{fase} comes from Greek \emph{phásis} 'an appearance' (EtymDB), the word for each new look of the moon.
\end{cousinweb}

\begin{grammarlens}[title={What you've built}]
One more word for telling, explaining and describing at B1.
\end{grammarlens}

\subsection*{Your turn}
\begin{itemize}
\raggedright
  \item \emph{Say it:} \emph{la fase}
  \item \emph{Say it:} \emph{la fase}, once more
  \item \emph{Say it:} a sentence of your own with \emph{la fase}
  \item \emph{Recall:} say the Spanish for weekly, then the Spanish for yearly, then say \emph{la fase} again
\end{itemize}

\begin{tcolorbox}[breakable,colback=teal!4,colframe=teal!35!black,title={Before you move on}]
What does la fase mean? (\textquotedblleft{}The phase, the stage\textquotedblright{}.) Say it once more.
\end{tcolorbox}
\section[el transcurso]{el transcurso — the passing, the span (of time)}

\label{lesson:ES-C541-transcurso}



\begin{quote}
Before the new one: say the Spanish for yearly, then the Spanish for the phase.
\end{quote}

\subsection*{You'll want to know: el transcurso}
\textbf{el transcurso} — \textquotedblleft{}the passing, the span (of time)\textquotedblright{}.

\begin{quote}
En el transcurso de la noche llegó más gente.
\end{quote}

\begin{quote}
\emph{As the night went on, more people arrived.}
\end{quote}

\emph{En el transcurso de} — during, over the span of.

\begin{cousinweb}
\emph{transcurso} is a learned borrowing from Latin \emph{trānscursus} 'a running through' (Lewis \& Short), built on \emph{trāns} 'across' and \emph{currere} 'to run'.
\end{cousinweb}

\begin{grammarlens}[title={What you've built}]
One more word for telling, explaining and describing at B1.
\end{grammarlens}

\subsection*{Your turn}
\begin{itemize}
\raggedright
  \item \emph{Say it:} \emph{el transcurso}
  \item \emph{Say it:} \emph{el transcurso}, once more
  \item \emph{Say it:} a sentence of your own with \emph{el transcurso}
  \item \emph{Recall:} say the Spanish for yearly, then the Spanish for the phase, then say \emph{el transcurso} again
\end{itemize}

\begin{tcolorbox}[breakable,colback=teal!4,colframe=teal!35!black,title={Before you move on}]
What does el transcurso mean? (\textquotedblleft{}The passing, the span (of time)\textquotedblright{}.) Say it once more.
\end{tcolorbox}
\section[la duración]{la duración — the length, the duration}

\label{lesson:ES-C541-duracion}



\begin{quote}
Before the new one: say the Spanish for the phase, then the Spanish for the passing.
\end{quote}

\subsection*{You'll want to know: la duración}
\textbf{la duración} — \textquotedblleft{}the length, the duration\textquotedblright{}.

\begin{quote}
La duración del viaje es de dos horas.
\end{quote}

\begin{quote}
\emph{The journey takes two hours.}
\end{quote}

\begin{cousinweb}
\emph{duración} is a learned borrowing from Medieval Latin \emph{dūrātiō} (EtymDB), built on \emph{dūrāre} 'to last', from \emph{dūrus} 'hard'.
\end{cousinweb}

\begin{grammarlens}[title={What you've built}]
One more word for telling, explaining and describing at B1.
\end{grammarlens}

\subsection*{Your turn}
\begin{itemize}
\raggedright
  \item \emph{Say it:} \emph{la duración}
  \item \emph{Say it:} \emph{la duración}, once more
  \item \emph{Say it:} a sentence of your own with \emph{la duración}
  \item \emph{Recall:} say the Spanish for the phase, then the Spanish for the passing, then say \emph{la duración} again
\end{itemize}

\begin{tcolorbox}[breakable,colback=teal!4,colframe=teal!35!black,title={Before you move on}]
What does la duración mean? (\textquotedblleft{}The length, the duration\textquotedblright{}.) Say it once more.
\end{tcolorbox}
\section[prolongarse]{prolongarse — to go on, to drag on}

\label{lesson:ES-C541-prolongarse}



\begin{quote}
Before the new one: say the Spanish for the passing, then the Spanish for the length.
\end{quote}

\subsection*{You'll want to know: prolongarse}
\textbf{prolongarse} — \textquotedblleft{}to go on, to drag on\textquotedblright{}.

\begin{quote}
La reunión se prolongó hasta la medianoche.
\end{quote}

\begin{quote}
\emph{The meeting went on until midnight.}
\end{quote}

It takes \emph{gu} before \emph{e}: \emph{se prolongue}. Without \emph{se}, \emph{prolongar} is to extend.

\begin{cousinweb}
\emph{prolongarse} is built on \emph{prolongar}, a learned borrowing from Late Latin \emph{prōlongāre} 'to lengthen' (EtymDB), built on \emph{prō} and \emph{longus} 'long' (Lewis \& Short).
\end{cousinweb}

\begin{grammarlens}[title={What you've built}]
One more word for telling, explaining and describing at B1.
\end{grammarlens}

\subsection*{Your turn}
\begin{itemize}
\raggedright
  \item \emph{Say it:} \emph{prolongarse}
  \item \emph{Say it:} \emph{prolongarse}, once more
  \item \emph{Say it:} a sentence of your own with \emph{prolongarse}
  \item \emph{Recall:} say the Spanish for the passing, then the Spanish for the length, then say \emph{prolongarse} again
\end{itemize}

\begin{tcolorbox}[breakable,colback=teal!4,colframe=teal!35!black,title={Before you move on}]
What does prolongarse mean? (\textquotedblleft{}To go on, to drag on\textquotedblright{}.) Say it once more.
\end{tcolorbox}
\section[transcurrir]{transcurrir — to go by, to pass (of time); to take place}

\label{lesson:ES-C541-transcurrir}



\begin{quote}
Before the new one: say the Spanish for the length, then the Spanish for to go on.
\end{quote}

\subsection*{You'll want to know: transcurrir}
\textbf{transcurrir} — \textquotedblleft{}to go by, to pass (of time); to take place\textquotedblright{}.

\begin{quote}
Transcurrieron dos horas sin noticias.
\end{quote}

\begin{quote}
\emph{Two hours went by without news.}
\end{quote}

A written word; the story \emph{transcurre en Sevilla} is set in Seville.

\begin{cousinweb}
\emph{transcurrir} is a learned borrowing from Latin \emph{trānscurrere} 'to run across' (EtymDB), built on \emph{trāns} 'across' and \emph{currere} 'to run'.
\end{cousinweb}

\begin{grammarlens}[title={What you've built}]
One more word for telling, explaining and describing at B1.
\end{grammarlens}

\subsection*{Your turn}
\begin{itemize}
\raggedright
  \item \emph{Say it:} \emph{transcurrir}
  \item \emph{Say it:} \emph{transcurrir}, once more
  \item \emph{Say it:} a sentence of your own with \emph{transcurrir}
  \item \emph{Recall:} say the Spanish for the length, then the Spanish for to go on, then say \emph{transcurrir} again
\end{itemize}

\begin{tcolorbox}[breakable,colback=teal!4,colframe=teal!35!black,title={Before you move on}]
What does transcurrir mean? (\textquotedblleft{}To go by, to pass (of time); to take place\textquotedblright{}.) Say it once more.
\end{tcolorbox}
